% ECONOMICS_PROBLEM: {"id": "C54-125", "title": "Robust structural counterfactuals", "jel_code": "C54", "source_id": "OP-0426"}
% !TeX program = xelatex
\documentclass[11pt,a4paper]{article}
\usepackage[margin=23mm]{geometry}
\usepackage{fontspec}
\setmainfont{Latin Modern Roman}
\usepackage{amsmath,amssymb,mathtools}
\usepackage{xcolor,enumitem,booktabs,longtable,array,xurl,hyperref}
\definecolor{muted}{HTML}{53636C}
\hypersetup{colorlinks=true,urlcolor=blue,linkcolor=blue}
\setlength{\parindent}{0pt}
\setlength{\parskip}{5pt}
\newcommand{\R}{\mathbb R}
\newcommand{\E}{\mathbb E}
\newcommand{\N}{\mathbb N}
\newcommand{\ind}{\mathbf 1}
\newcommand{\Var}{\operatorname{Var}}
\newcommand{\Cov}{\operatorname{Cov}}
\newcommand{\supp}{\operatorname{supp}}
\DeclareMathOperator*{\argmin}{arg\,min}
\DeclareMathOperator*{\argmax}{arg\,max}
\newcommand{\entrymeta}[3]{{\sffamily\small\color{muted}\textbf{\texttt{#1}}\quad|\quad #2\par #3\par}}
\newcommand{\Problemref}[1]{\texttt{#1}}
\newcommand{\JELref}[2]{\texttt{#2} (JEL \texttt{#1})}
\begin{document}
\section*{Robust structural counterfactuals}
\textbf{Problem ID:} \texttt{C54-125}\quad \textbf{JEL:} \texttt{C54}\par
% BEGIN ECONOMICS STATEMENT
\label{problem:OP-0426}
{\sffamily\small\color{muted}Original subject group: Econometrics, causal inference and measurement\par}
\label{definitions:problem:30007708}
\textbf{Audit classification:} Background; proposed extension. Definition6 and Proposition3 explicitly rule out weak identification. The proposed simultaneous weak-IV and total-variation counterfactual frontier is not stated as open by this source; coverage alone would be vacuous.
\subsection*{Definitions and precise research target}
Fix a structural model family \(M=(\theta,F_U)\) mapping observed covariates \(X\), endogenous choices \(D\) and latent variables \(U\) into outcomes through a known function \(g(X,D,U;\theta)\); \(\theta\in\Theta\), and \(F_U\) is a latent distribution. Observed instruments \(Z\) satisfy stated exogeneity restrictions in the exact model. Specified choice equations and exogenous-variable laws complete the family; let \(P_M\) be its implied observed-data distribution. For actual observed law \(P\), tolerance \(\delta\in[0,1]\), and total-variation distance \(d_{\rm TV}(P,P_M)=\sup_B|P(B)-P_M(B)|\) over measurable events \(B\), define \(\mathcal M_\delta(P)=\{M:d_{\rm TV}(P,P_M)\leq\delta\}\). For a fully specified policy \(q\), let \(\psi(M,q)\) be the resulting average welfare change after the model's equilibrium choices adjust under a declared selection rule. Assume a known finite bound \(B_0\) with \(|\psi(M,q)|\leq B_0\) for all admitted models. Construct confidence sets covering \(\{\psi(M,q):M\in\mathcal M_\delta(P)\}\), uniformly over independent-sample data laws as instrument relevance approaches zero. For a stated scalar linear-IV moment, weakness is measured by the conditional covariance of \(D,Z\) given \(X\); state the multivariate analogue if used. Compare expected set width across tolerances, for sample size \(n\) and significance level \(\alpha\in(0,1)\), and identify validation data justifying smaller \(\delta\). Existing papers establish robustness in important settings; this total-variation counterfactual frontier under simultaneous weakness and misspecification is a proposed specification, not a verified canonical open problem.

\textbf{Definitions, assumptions and mathematical qualifications.} A complete structural model includes choice equations, exogenous-variable and latent laws, equilibrium sets and a selected policy outcome; its observation map is \(M\mapsto P_M\). Require \(\mathcal M_\delta(P)\ne\varnothing\). A confidence region has identified-set coverage if \(\inf_{P\in\mathcal P}\Pr_{P^n}(\{\psi(M,q):M\in\mathcal M_\delta(P)\}\subseteq C_n)\geq1-\alpha\). The fixed region \([-B_0,B_0]\) trivially satisfies coverage; the substantive task is to attain useful expected width or minimax excess length. For a scalar IV equation use standardized residualized first-stage covariance, or in multivariate moments the smallest singular value of the normalized relevant Jacobian, to define weakness. Observational total variation does not bound counterfactual distance: observationally equivalent models can predict different policies even at \(\delta=0\).
\subsection*{Variants and assumption changes}
\begin{enumerate}
\item \textbf{Sourced strong-identification benchmark.} The author draft strong-exclusion robustness result imposes strong identification; apply it only under its declared regularity and instrument conditions.
\item \textbf{Editorial weak/misspecified extension.} Specify a bounded model family, counterfactual regularity and a meaningful width criterion; seek matching upper/lower minimax bounds as the first-stage strength tends to zero while \(\delta\) remains positive.
\end{enumerate}
\subsection*{References and source audit}
\textbf{Checked source.} January 2025 author draft, Section 4, Definition 6 and Proposition 3, printed pp. 26--27. The robustness result maintains strong identification and explicitly excludes weak identification. Full author text checked; initial NBER working paper issued 2023, DOI 10.3386/w31799; published QJE 140(3) (2025), pp. 1801--1855, DOI 10.1093/qje/qjaf018. \url{https://web.stanford.edu/~gentzkow/research/includediv.pdf}.
\begin{enumerate}\raggedright
\item\hypertarget{definitions:source:30007708:1}{} Isaiah Andrews; Nano Barahona; Matthew Gentzkow; Ashesh Rambachan; Jesse M. Shapiro. 2025. \emph{Structural Estimation Under Misspecification: Theory and Implications for Practice}. January 2025 author draft, Section 4, Definition 6 and Proposition 3, printed pp. 26--27.
\url{https://web.stanford.edu/~gentzkow/research/includediv.pdf}.
DOI: \url{https://doi.org/10.3386/w31799}.
\end{enumerate}

\subsection*{Common conventions: Source mathematical conventions}

These conventions apply to each entry unless explicitly replaced.

\textbf{Models and identification.} A primitive model \(m\) specifies all probability laws, feasible choices, information, equilibrium and measurement rules used by its target. The admissible class \(\mathcal M\) is fixed independently of outcomes. If \(P_O(m)\) is its observable probability law and \(\tau(m)\) the target, its identified set at observable law \(P\) is
\[
 \mathcal I_\tau(P)=\{\tau(m):m\in\mathcal M,\ P_O(m)=P\}.
\]
Point identification means that this set is a singleton. An empty set signals incompatibility of data and assumptions. Coordinate bounds need not describe the joint set. Bounds are sharp only if every retained target is attainable or arbitrarily approachable by compatible models. Finite bounds require explicit restrictions on unobserved quantities; unrestricted confounding, hidden wealth, extrapolation or arbitrary equilibrium selection can yield uninformative sets.

\textbf{Causal interventions.} A potential outcome \(Y_i(p)\) is the outcome under a complete policy regime \(p\), including timing, financing, information and market spillovers. The target population and units are fixed before treatment. With interference, use the complete assignment vector or a justified exposure mapping; individual no-interference notation alone cannot identify an economy-wide intervention. A difference of mean potential outcomes does not require the cross-world joint distribution. Conditioning on post-treatment migration, survival, enrollment or employment changes the estimand and needs separate justification.

\textbf{Statistical statements.} An estimator, confidence set, forecast or test specifies a sampling law, model class and loss/coverage criterion. Uniform \((1-\alpha)\)-coverage means \(\inf_{m\in\mathcal M}\Pr_m\{\tau(m)\in C_n\}\geq1-\alpha\), or an explicitly asymptotic version. The whole parameter space trivially has coverage, so informativeness requires a stated width, risk or power objective. A forecasting improvement is relative to a named benchmark, information set and evaluation population.

\textbf{Equilibrium and optimization.} An equilibrium comprises feasible plans, optimality, beliefs where needed, budget constraints and all market/resource clearing. The primitives or a stated benchmark define each correspondence \(\mathcal E(p;m)\). Nonemptiness is assumed or proved, never inferred just from compact policy sets. A selected-equilibrium target uses a declared selection rule; a robust target quantifies over all permitted equilibria. Use suprema or infima unless attainment is established. An \(\varepsilon\)-optimal policy is feasible and within \(\varepsilon\) of the finite optimum value. Report an empty feasible set separately.

\textbf{Welfare and accounting.} Normative utility, interpersonal normalization, social weights, numeraire, discounting and the use of public receipts are primitives. Behavioral utility need not coincide with the welfare criterion. Household welfare includes ownership income and rents once; adding profits again to compensating variation generally double-counts them. A monetary equivalent is defined by an explicit transfer and price convention, with monotonicity and range conditions. Average effects, mechanism contributions, transfers and welfare losses are different targets. Infinite sums require convergence and dynamic feasible plans require terminal or no-Ponzi conditions.

\textbf{Source versions.} A working-paper date, revision date and journal publication year are distinguished. A DOI for a journal version is not automatically the DOI of an earlier manuscript. Locators refer to the stated version and printed pages where specified. A metadata-only check does not verify a claim about a full-text conclusion. Every entry's source subsection narrows the provenance claim accordingly.
% END ECONOMICS STATEMENT
\end{document}
